\section{Approximation for power-weight and reversed-weight amplitudes}\label{sec:companion}

We prove the approximation with the power-weight norm $\pwNorm a$ and transfer it to the reversed-weight norm $\swNorm a$, which gives Theorem~\ref{thm:dd14}. These arguments exhibit the common remainder estimate and the balance between analytic smallness and weighted transport growth. Section~\ref{sec:printed} then treats the original norm, which requires additional critical-weight estimates.
Throughout, $M=\Gamma\backslash\Hb$ and $d=d(x_0,\cdot)$.

\medskip\noindent\emph{The common structure.}
The proof of Theorem~\ref{thm:dd-appx} and that of Proposition~\ref{pc:core} in Section~\ref{subsec:pc-proof} both bound the
integrand $\norm{\mathcal V_{s_1}a_{\mathbf s}}_{L^2(M)}$ of Proposition~\ref{prop:dd-id}(ii), with $\varepsilon=(\alpha\hbar)^{1/2}/2$
and $\xi=(\alpha/\hbar)^{1/2}$ ($\sigma$ in place of $\alpha$ in Proposition~\ref{pc:core}), by a mollified part and a defect
part as described in Section~\ref{subsec:method}, and then integrate over the simplex with Lemma~\ref{lem:collect}.
Table~\ref{tab:compare} lists where they differ.

\begin{center}
\small
\captionof{table}{The power-weight proof and the original-norm proof. The operator $\Lambda_m$ is defined in
Section~\ref{subsec:pc-weights}, and the weight $\omega$ in Section~\ref{sec:setup}.}\label{tab:compare}
\begin{tabular}{@{}>{\raggedright\arraybackslash}p{0.18\linewidth}>{\raggedright\arraybackslash}p{0.34\linewidth}>{\raggedright\arraybackslash}p{0.40\linewidth}@{}}
\hline
Stage & Power-weight proof (Section~\ref{subsec:dd-thm}) & Original-norm proof (Section~\ref{sec:printed})\\
\hline
Data & $a$ with $\pwNorm a<\infty$ & $a=\ue^{-\sigma\Lambda_m}\tilde f$ with $\omega\tilde f\in L^2$ (Proposition~\ref{os:data})\\[4pt]
Weighted bounds & H\"older interpolation between $L^2$ and the weight $X_0^\beta$ & interpolation shell by shell between $L^2$ and the weight $X_0$ (Lemma~\ref{os:transfer})\\[4pt]
Periodization & weight $X_0^\lambda$, $\lambda>\tfrac12$ (Lemma~\ref{lem:t2}) & weight $\omega=X_0^{1/2}(1+\tilde\psi/\log2)^{3/4}$ (Proposition~\ref{os:p1})\\
\hline
\end{tabular}
\end{center}

\subsection{Approximation with the power weight \texorpdfstring{$X_0^\beta$}{X0}}\label{subsec:dd-thm}

\begin{theorem}[approximation in the power-weight class]\label{thm:dd-appx}
Fix $\beta>\tfrac12$, $\kappa>0$ and $A>0$. The statement of Theorem~\ref{pc:main} holds, with constants $\delta$ and $C$ of its
own, when $\pcNorm a$ is replaced by $\pwNorm a=\norm{\ue^{\alpha\Lambda}X_0^\beta a}$ and the series $P^{\mathrm{cr}}_\Gamma$ by
$P_\Gamma(\beta+\tfrac12)$.
\end{theorem}

Steps~0--3 below prove the estimate for $0<\alpha\le\min(A,1)$, where the estimates of Section~\ref{sec:analytic} apply; the
last paragraph extends it to $1<\alpha\le A$.

\begin{proof}
\medskip\noindent\emph{Step 0 (reductions and parameters).} By Lemma~\ref{lem:base} we may take $x_0=i$ and $b=\infty$: the group enters the estimates below only
through the Poincar\'e series bound at the base point. By Lemma~\ref{lem:lift}, $a$ is $C^\infty$. We write $\varphi=\varphi_\infty$ and
\[
 E:=\Delta^N\bigl[\cU(t)[a\ue^{i\varphi/\hbar}]_\Gamma-u^{(K)}(t)_\Gamma\bigr].
\]

Let $c_M$ be the constant of Lemma~\ref{lem:mw} and $c_A$ that of
Proposition~\ref{prop:a0}. Put
\[
r_\alpha:=c_M\alpha,\qquad r_1:=\tfrac34r_\alpha,\qquad r_2:=\tfrac12r_\alpha,\qquad\varepsilon:=(\alpha\hbar)^{1/2}/2,\qquad \xi:=(\alpha/\hbar)^{1/2},
\]
so that $\varepsilon\le\tfrac12$, $\varepsilon/\hbar=\xi/2$ and $\alpha/\varepsilon=2\xi$. Fix $\lambda:=(2\beta+1)/4\in(1/2,\beta)$ and put $\theta:=\lambda/\beta\in(0,1)$.

Use Proposition~\ref{prop:dd-id}(ii) with $\varpi=X_0^\lambda$; every word applied to $a$ lies in $L^2_\varpi$ by Lemma~\ref{lem:sc}, since $X_0^\lambda\le X_0^\beta$. For each $\mathbf s=(s_1,\dots,s_{K+1})\in\Sigma_{K+1}(t)$ let
\[
g(t'):=V_\varepsilon(t',s_1)\,a_{\mathbf s},\quad0\le t'\le s_1,\qquad w:=g(0)=V_\varepsilon(s_1)^*a_{\mathbf s},
\]
and write, since $g(s_1)=a_{\mathbf s}$,
\[
 \mathcal V_{s_1}a_{\mathbf s}=\mathcal V_0w+\bigl(\mathcal V_{s_1}g(s_1)-\mathcal V_0g(0)\bigr).
\]
The first term is bounded by weighted propagation of the mollified flow (Step~1). For the second, the Duhamel formula on $M$ (Lemma~\ref{lem:surf-duh})
expresses the difference through the mollifier defect, whose analytic smallness compensates for the transport growth
after interpolation (Step~2).

\medskip\noindent\emph{Step 1 (mollified part).} Lemma~\ref{lem:t2} at $t=0$, and Proposition~\ref{prop:ws} at the weight $X_0^\lambda$ (with $\hbar s_1\le\varepsilon$), give
\[
\norm{\mathcal V_0w}=\norm{[\ue^{i\varphi/\hbar}V_\varepsilon(s_1)^*a_{\mathbf s}]_\Gamma}\le C\norm{X_0^\lambda V_\varepsilon(0,s_1)a_{\mathbf s}}\le C\ue^{C_{\mathrm{fl},\lambda}/2}\norm{X_0^\lambda a_{\mathbf s}}\le C\norm{X_0^\beta a_{\mathbf s}},
\]
using $X_0^\lambda\le X_0^\beta$ (since $X_0\ge1$ and $\lambda\le\beta$).
By Lemma~\ref{lem:wb}(i) with $n=0$:
\[
\norm{\mathcal V_0w}\le C\,C_3^{N+K+1}(2N)!\,(2K+2)!\,\alpha^{-(2K+2)}(\alpha\hbar)^{-2N}\pwNorm a.
\]


\medskip\noindent\emph{Step 2 (defect part).} By Lemma~\ref{lem:wb}(i) with $n\le2$, $X_0^\beta Za_{\mathbf s}\in L^2$ for every $Z$ in $\mathcal S=\{1,X,Y,X^2,Y^2\}$. Thus Lemma~\ref{lem:path}(c), with the exponent
$\lambda$, shows that $g$ satisfies the hypotheses of Lemma~\ref{lem:surf-duh} for $\varpi=X_0^\lambda$. Since $g'=c_\hbar\,\Delta_\varepsilon(t')g$,
Duhamel on $M$ gives $\mathcal V_{s_1}g(s_1)-\mathcal V_0g(0)=-c_\hbar\int_0^{s_1}\mathcal V_{t'}\mathfrak d(t')\,dt'$, with
\[
\mathfrak d(t'):=(\Delta(t')-\Delta_\varepsilon(t'))\,g(t').
\]
By the unitarity of $\cU(t')$ on $M$ and Lemma~\ref{lem:t2} (exponent $\lambda$),
\[\norm{\mathcal V_{t'}\mathfrak d(t')}=\norm{[\ue^{i\varphi/\hbar}S(t')\mathfrak d(t')]_\Gamma}\le C\ue^{\lambda t'}\norm{X_0^\lambda \mathfrak d(t')}.\]
H\"older's inequality gives $\norm{X_0^\lambda\mathfrak d}\le\norm{X_0^\beta\mathfrak d}^\theta\norm{\mathfrak d}^{1-\theta}$, since $X_0^\lambda=(X_0^\beta)^\theta$.
We estimate the two factors separately.

\emph{Weighted factor.} By Lemma~\ref{lem:path}(b) with the exponent $\beta$, and Lemma~\ref{lem:wb}(i) with $n\le2$ (five words; the bound increases with $n$),
\begin{align*}
\norm{X_0^\beta \mathfrak d(t')}&\le C\varepsilon^{-3}\sum_{Z\in\mathcal S}\norm{X_0^\beta Za_{\mathbf s}}\\
&\le C\varepsilon^{-3}\,C_3^{N+K+3}(2N)!\,(2K+4)!\,\alpha^{-(2K+4)}(\alpha\hbar)^{-2N}\pwNorm a.
\end{align*}
\emph{Unweighted factor.} Lemma~\ref{lem:mw} gives $\norm{\ue^{r_\alpha\Lambda}a}\le C_*B_\beta\pwNorm a$.
Lemma~\ref{lem:wb}(ii), applied with $r_\alpha$ in the role of $\alpha$ and the radii $r_1<r_\alpha$ (gap $r_\alpha/4$), gives
$\norm{\ue^{r_1\Lambda}a_{\mathbf s}}\le C_3^{N+K+1}(2N)!(2K+2)!\,r_\alpha^{-(2K+2)}(r_\alpha\hbar)^{-2N}\norm{\ue^{r_\alpha\Lambda}a}$.
Proposition~\ref{prop:a0}, from the radius $r_1$ to $r_2$ (so $(r_1-r_2)/r_1=\tfrac13$) and for $|t'-s_1|\le t\le\tfrac13c_A\varepsilon/\hbar$, gives
$\norm{\ue^{r_2\Lambda}g(t')}\le\norm{\ue^{r_1\Lambda}a_{\mathbf s}}$.
Lemma~\ref{lem:unw} with $r=r_2$ gives $\norm{\mathfrak d(t')}\le C_4r_2^{-4}\ue^{-r_2/(8\varepsilon)}\norm{\ue^{r_2\Lambda}g(t')}$, and $r_2/(8\varepsilon)=c_M\xi/8$.

\emph{Combining the factors.} Let $F_n$ be the numbers of Lemma~\ref{lem:collect} for $G:=C_3/c_M^2$, $\sigma:=\alpha$ and
$\mathcal N:=\pwNorm a$. Since $c_M\le1\le C_3$ and $\alpha\le1$, the weighted factor is at most $C\varepsilon^{-3}F_2$ and the unweighted one at
most $C\alpha^{-4}\ue^{-c_M\xi/8}F_2$. After the H\"older split, using $\varepsilon,\alpha\le1$, the defect part is at most
\[
C\,\hbar t\,\ue^{\lambda t}\,\varepsilon^{-3}\alpha^{-4}\ue^{-c\xi}F_2,\qquad c:=(1-\theta)c_M/8 .
\]

\medskip\noindent\emph{Step 3 (collecting).} The bound of Step~1 is at most $CF_0$. So for every $\mathbf s\in\Sigma_{K+1}(t)$,
\[
\norm{\mathcal V_{s_1}a_{\mathbf s}}\le B\bigl(F_0+\hbar t\,\ue^{\lambda t}\varepsilon^{-3}\alpha^{-4}\ue^{-c\xi}F_2\bigr),
\]
with $B$ depending only on $\beta$, $\kappa$ and $P$. By Proposition~\ref{prop:dd-id}(ii), the error $\norm E$ is at most $(\hbar/2)^{K+1}$
times the volume $t^{K+1}/(K+1)!$ of $\Sigma_{K+1}(t)$ times this bound, and Lemma~\ref{lem:collect}, with the exponent $\lambda$ of Step~0 as its parameter $\lambda$, gives
\[
\norm E\le C^{N+K+1}N^{2N}K!\,(\alpha\hbar)^{-2N}\Bigl(\frac{\hbar t}{\alpha^2}\Bigr)^{K+1}\pwNorm a
\]
for $t\le\delta\xi$ with $\delta\le1$ and $2\lambda\delta\le c$.

\medskip\noindent\emph{Time windows used.} Weighted propagation and Lemma~\ref{lem:path} require $\hbar s_1\le\varepsilon$, that is $t\le\xi/2$.
Analytic propagation requires $t\le\tfrac13c_A\varepsilon/\hbar=c_A\xi/6$.
Lemma~\ref{lem:collect} requires $2\lambda\delta\le c$, that is $\delta\le(1-\theta)c_M/(16\lambda)$.

So $\delta:=\min\bigl(\tfrac12,\;c_A/6,\;(1-\theta)c_M/(16\lambda)\bigr)$ works.

\medskip\noindent\emph{Dependence on $\Gamma$ and $x_0$.} The reduction of Step~0 preserves $P_\Gamma(\beta+\tfrac12)$
(Lemma~\ref{lem:base}(iii)). The group enters only through Lemma~\ref{lem:t2}, used in Steps~1 and~2 with exponent $\lambda$, whose
constant $C(\Gamma,\lambda)=2^\lambda P_\Gamma(2\lambda)^{1/2}$ depends on $\Gamma$ only through the Poincar\'e series bound $P_\Gamma(2\lambda)$, $2\lambda=\beta+\tfrac12$.
All other inputs are estimates on $\Hb$. So $\delta$ does not depend on $\Gamma$ or $x_0$, and $C$ depends on them only through
$P_\Gamma(\beta+\tfrac12)$.

\medskip\noindent\emph{Radii $1<\alpha\le A$.} The power-weight norm increases with the radius, since
$\ue^{\Lambda}=\ue^{-(\alpha-1)\Lambda}\ue^{\alpha\Lambda}$ and $\ue^{-\tau\Lambda}$ is a contraction for $\tau\ge0$. Apply the case already proved
at the radius $1$ with $1$ in place of $\kappa$ (then $1\ge\hbar$), and let $\bar C,\bar\delta$ be its constants. For
$t\le\bar\delta A^{-1/2}(\alpha/\hbar)^{1/2}$ we have $t\le\bar\delta\hbar^{-1/2}$, and
\[
\hbar^{-2N}(\hbar t)^{K+1}\le A^{2N+2K+2}(\alpha\hbar)^{-2N}(\hbar t/\alpha^2)^{K+1}.
\]
So the bound holds for $1<\alpha\le A$ with
$\bar\delta A^{-1/2}$ and $A^2\bar C$; take the smaller of the two values of $\delta$ and the larger of the two values of $C$. Both
cases have the stated dependence on $\Gamma$ and $x_0$.
\end{proof}


\subsection{Proof of Theorem~\texorpdfstring{\ref{thm:dd14}}{1.2}}\label{sec:swapped}
We now transfer the power-weight estimate to the norm $\swNorm a=\norm{\ue^{\alpha\Lambda}\ue^{\beta\langle d\rangle}a}$, with spatial
decay imposed on the amplitude, first for radii $\alpha\le\tfrac12$ and then, by the monotonicity of this norm in the radius, for
every bounded range of radii. By Lemma~\ref{lem:m2w}, a finite norm $\swNorm a$ at radius $\alpha$ gives a finite norm with the
weight $X_0^\beta=\ue^{\beta\tilde\psi}$ in place of $\ue^{\beta\langle d\rangle}$, at the radius $c_0\alpha$ for a fixed
$c_0=c_0(\beta)\in(0,1]$.

\begin{proof}[Proof of Theorem~\ref{thm:dd14}]
\emph{Radii $0<\alpha\le\tfrac12$.} Let $c_0$ and $C'$ be the constants $c_0$ and $C$ of Lemma~\ref{lem:m2w}. If $\alpha\le\tfrac12$ and
$\alpha^3\ge\kappa\hbar$, then $\bar\alpha:=c_0\alpha\le c_0/2$, $\bar\alpha^3\ge c_0^3\kappa\hbar$ and
$\norm{\ue^{\bar\alpha\Lambda}X_0^\beta a}\le C'\swNorm a$. Apply Theorem~\ref{thm:dd-appx} with $c_0/2$, $c_0^3\kappa$ and $\bar\alpha$
in place of $A$, $\kappa$ and $\alpha$, and let $C_{\rm pw},\delta_{\rm pw}$ be its constants. Its
window $t\le\delta_{\rm pw}(\bar\alpha/\hbar)^{1/2}$ is the window $t\le\delta_{\rm s}(\alpha/\hbar)^{1/2}$ with $\delta_{\rm s}:=c_0^{1/2}\delta_{\rm pw}$, and
$(\bar\alpha\hbar)^{-2N}(\hbar t/\bar\alpha^2)^{K+1}=c_0^{-2N-2K-2}(\alpha\hbar)^{-2N}(\hbar t/\alpha^2)^{K+1}$. So the bound holds for
$\alpha\le\tfrac12$ with the constants $\delta_{\rm s}$ and $C_{\rm s}:=C_{\rm pw}\,c_0^{-2}\max(1,C')$.

\emph{Radii $\tfrac12<\alpha\le A$.} The reversed-weight norm increases with the radius: $\ue^{\Lambda/2}=\ue^{-(\alpha-1/2)\Lambda}\ue^{\alpha\Lambda}$,
and $\ue^{-\tau\Lambda}$ is a contraction for $\tau\ge0$, so the norm of $a$ at the radius $\tfrac12$ is at most $\swNorm a$. Apply the first
case at the radius $\tfrac12$ with $\tfrac18$ in place of $\kappa$ (then $(\tfrac12)^3\ge\tfrac18\hbar$ for $\hbar\le1$), and let $C_{1/2},\delta_{1/2}$ be its
constants. For $t\le\delta_{1/2}(2A)^{-1/2}(\alpha/\hbar)^{1/2}$ we have $t\le\delta_{1/2}(1/(2\hbar))^{1/2}$, and
$(\tfrac12\hbar)^{-2N}(4\hbar t)^{K+1}\le(2A)^{2N+2K+2}(\alpha\hbar)^{-2N}(\hbar t/\alpha^2)^{K+1}$. So the bound holds for
$\tfrac12<\alpha\le A$ with $\delta_{1/2}(2A)^{-1/2}$ and $(2A)^2C_{1/2}$. Take $\delta:=\min(\delta_{\rm s},\delta_{1/2}(2A)^{-1/2})$ and
$C:=\max(C_{\rm s},(2A)^2C_{1/2})$.

In both cases the group and the base point enter only through Theorem~\ref{thm:dd-appx}, whose $\delta_{\rm pw}$ does not depend
on them and whose $C_{\rm pw}$ depends on them only through $P_\Gamma(\beta+\tfrac12)$; Lemma~\ref{lem:m2w} holds for every base
point with the same constants. So the same holds for $\delta$ and $C$.
\end{proof}
